% SEGMENT OLP-0546-S01
% Part: set-theory
% Chapter: z
% Section: arbintersections

\documentclass[../../../include/open-logic-section]{subfiles}

\begin{document}

\olfileid{sth}{z}{arbintersections}
\olsection{பின்னிணைப்பு: மூடல், பண்புவழிக் கணமாக்கல், வெட்டு}

\olref[sth][z][milestone]{sec} இல், \olref[sfr][][]{part} இல்
முறைசாராமல் செய்தவற்றைத் திரும்பிப் பார்த்து, அவற்றை $\Zminus$ இல்
செய்ய முடியும் எனச் சரிபார்க்குமாறு பரிந்துரைத்தோம். அந்த அறிவுரையை நீங்கள்
பின்பற்றியிருந்தால், டெடெகிண்டின் \emph{மூடல்களைக்} கையாளும்போது
\emph{வெட்டை} பயன்படுத்தியது என்ற \emph{ஒரு} புள்ளி உங்களைத்
தடுமாறச் செய்திருக்கலாம்.

\olref[sfr][infinite][dedekind]{Closure} இலிருந்து பின்வருவதை நினைவுகூர்க:
\begin{align*}
  \closureofunder{f}{o} & = \bigcap\Setabs{X}{o \in X \text{ மற்றும் $X$, $f$-மூடப்பட்டது}}.
\intertext{இதன் பொதுவடிவம் பின்வரும் வடிவிலான ஒரு வரையறை:}
  C & = \bigcap\Setabs{X}{\phi(X)}.
\end{align*}
ஆனால் இது எச்சரிக்கை மணியை ஒலிக்கச் செய்ய வேண்டும்: முறைசாராப்
பண்புவழிக் கணமாக்கல் தோல்வியடைவதால், $\Setabs{X}{\phi(X)}$ உள்ளது
என்பதற்கு எந்த உறுதியும் இல்லை. அப்படியானால், இத்தகைய வரையறைகள்
\emph{ஏமாற்றுவதாக} இருப்பதுபோல் ஆபத்தான விதத்தில் தோன்றுகிறது.

நல்வாய்ப்பாக, அவை ஏமாற்றுவதில்லை; அல்லது, அவை இப்போதுள்ள வடிவில்
\emph{ஏமாற்றுகின்றன} என்றாலும், அவற்றை முறையாக ஏற்றுக்கொள்ளத்தக்கவையாக்க
நேர்மையான உழைப்பை மேற்கொள்ளலாம். எந்த $A \neq \emptyset$ க்கும்
$\bigcap A$ உள்ளது என்று விளக்கியபோது, அந்த உழைப்பை
\olref[sth][z][sep]{prop:intersectionsexist} இல் முன்னறிவித்தோம்.
ஆனால் அதனை இப்போது வெளிப்படையாகக் கூறுவோம்.

உறுப்புசார் சமத்துவம் தரப்பட்டிருக்க, $C$ ஐ
$\bigcap\Setabs{X}{\phi(X)}$ என வரையறுக்க முயலும்போது, பின்வருவதை
நிறைவேற்றும் ஒரு பொருள் $C$ ஐ மட்டுமே உண்மையில் கோருகிறோம்:
\begin{equation}\ollabel{bicondelimarbintersection}
	\forall x(x \in C \liff \forall X(\phi(X) \lif x \in X))\tag{*}
\end{equation}
இப்போது, $\phi(S)$ ஆக அமையும் \emph{ஏதோ ஒரு} கணம் $S$ உள்ளது என்க.
அப்போது \olref{bicondelimarbintersection} ஐப் பெற, பின்வருமாறு
\emph{பிரிப்பைப்} பயன்படுத்தி $C$ ஐ நேரடியாக வரையறுக்கலாம்:
\[
	C = \Setabs{x \in S}{\forall X(\phi(X) \lif x \in X)}.
\]
இந்த வரையறை விரும்பியவாறு \olref{bicondelimarbintersection} ஐத் தருகிறது
என்பதைச் சரிபார்ப்பதைப் பயிற்சியாக விடுகிறோம்.
%  $x$ ஐ நிலைப்படுத்துக. முதலில், (மறு)வரையறையின்படி $x \in C$ என்க; அப்போது $x
%  \in S$ மற்றும் $\forall X(\phi(X) \lif x \in X)$ ஆகிய இரண்டும் மெய்; ஆனால்
%  $\phi(S)$ என்பதால், இது $\forall X(\phi(X)
%  \lif x \in X)$ என்ற நிபந்தனையாகச் சுருங்குகிறது. மறுதிசையில், $\forall X(\phi(X) \lif
%  x \in X)$ என்க; அப்போது $\phi(S)$ என்பதால் $x \in S$ உம் கிடைக்கிறது; ஆகவே
%  (மறு)வரையறையின்படி $x \in C$.
% SOURCE CORRECTION TA-STH-009: restore uppercase C in the inactive explanatory proof.
இந்தப் பொது உத்தி, வெட்டுகளை வரையறுப்பதில் முறைசாராப் பண்புவழிக்
கணமாக்கல் பயன்படுத்தப்படுவது போல் தோன்றும் எந்த நேர்வையும் தவிர்க்க
உதவும். இந்தச் சிந்தனையைத் தொடங்கிய குறிப்பிட்ட நேர்வான
$\closureofunder{f}{o}$ இல் இது எவ்வாறு செயல்படும் என்பது இதோ.
$o \in \ran{f}\cup\{o\}$ என்றும் $\ran{f} \cup \{o\}$ என்பது
$f$-மூடப்பட்டது என்றும் குறிப்பிட்டு
\olref[sfr][infinite][dedekind]{closureproperties} இன் நிறுவலைத்
தொடங்கினோம். ஆகவே நாம் விரும்புவதைப் பின்வருமாறு வரையறுக்கலாம்:
\[
	\closureofunder{f}{o} = \Setabs{x \in \ran{f} \cup \{o\}}{(\forall X \ni o)(X \text{ ஆனது $f$-மூடப்பட்டது} \lif x \in X)}.
\]

\end{document}
